\section*{Chapter \ref{ch:ll:measuring-latency} --- Measuring Latency}

\begin{solution}{exo:ll:measuring-latency:1}
$(3\,020 - 45)/3.02 = \qty{985.1}{\nano\second}$.
\end{solution}

\begin{solution}{exo:ll:measuring-latency:2}
1\,000 has bit length 10, so the shift is 2 and the top bits are $\lfloor 1\,000/4 \rfloor = 250$: index $256 + 128 +
(250-128) = 506$, covering 1\,000 to 1\,003.
\end{solution}

\begin{solution}{exo:ll:measuring-latency:3}
$1\,024 + 53 \times 512 + 512 = 28\,672$ counters, with a relative precision of $2^{-9} \approx 0.195\%$.
\end{solution}

\begin{solution}{exo:ll:measuring-latency:4}
7, 5 and \qty{3}{\milli\second}: the values $9 - 2k$ while they exceed the \qty{2}{\milli\second} interval.
\end{solution}

\begin{solution}{exo:ll:measuring-latency:5}
The correction assumes that each missing request would have been answered as soon as the freeze ended. In fact the ten
requests scheduled during the freeze queue behind each other and are served one after another, so those sent later in
the freeze wait for the earlier ones as well; the open-loop tester measures that queue, the correction does not.
\end{solution}

\begin{solution}{exo:ll:measuring-latency:6}
It records only their processing time, as if they had arrived when the loop reached them: nothing of the up to
\qty{2}{\milli\second} they waited in the socket. It should measure from the packet's arrival: a \omterm{def:ll:measuring-latency:hwts}{hardware timestamp} from
the card, or at least the kernel's receive timestamp (\texttt{SO\_TIMESTAMPNS}).
\end{solution}

\begin{solution}{exo:ll:measuring-latency:7}
\texttt{ll\_measure.spectrum(stall=50\_000, period=5e6)}: closed loop \qty{0.16}{\milli\second}, corrected
\qty{1.07}{\milli\second}, open loop \qty{5.77}{\milli\second}. The closed-loop tester takes one sample per freeze, now
20 in 99\,000 (0.02\%), far below 1\%: its p99 does not see the freezes at all, while the users queued behind each
\qty{50}{\milli\second} freeze do.
\end{solution}

\begin{solution}{exo:ll:measuring-latency:8}
Three flaws: one closed-loop client (\omterm{def:ll:measuring-latency:co}{coordinated omission} hides every stall); linear buckets of \qty{100}{\nano\second}
up to \qty{100}{\micro\second} cannot hold the tail at all, so any value above the top bucket is lost or clamped; and a
million samples give only 100 above p99.99, too few to state it without a spread. Use an open-loop generator at the
production rate, a log-linear histogram, and repeated runs.
\end{solution}

\begin{solution}{pb:ll:measuring-latency:1}
\begin{enumerate}
\item The open-loop tester sends 100\,000; the closed-loop one 99\,000, because each freeze swallows the ten
milliseconds during which it waits for one reply.
\item 100 of 99\,000: 0.10\%.
\item The ten requests scheduled during each freeze, 1\% of all, plus those queued behind them just after it.
\item Nine per freeze (\qty{9.1}{\milli\second}, \qty{8.1}{\milli\second}, \dots, \qty{1.1}{\milli\second}): 900, so the
corrected histogram holds 99\,900 samples.
\item \qty{0.16}{\milli\second} (closed loop), \qty{1.06}{\milli\second} (corrected) and \qty{1.87}{\milli\second} (open
loop).
\item At the median every tester sees the ordinary \qty{100}{\micro\second} reply; beyond p99.9 all three see the freeze
itself, about \qty{10.1}{\milli\second}.
\item The correction assumes that the missing requests would have been served at the end of the freeze; the queue that
builds meanwhile makes them wait longer.
\item The open-loop figure: users send on their own schedule.
\item $2^{-7} \approx 0.8\%$; irrelevant here, where the testers disagree by factors of ten.
\item 7\,424.
\item Mechanically yes, by adding counts; but merging the output of different testers mixes different questions. Merge
histograms of the same tester across runs or hosts.
\item Every freeze sample, at about \qty{10.1}{\milli\second}, would fall outside the top bucket and be clamped or lost:
the report would have no tail.
\item \textbf{Named result.} The closed-loop tester reports a p99 of \qty{0.16}{\milli\second} raw and
\qty{1.06}{\milli\second} corrected; the open-loop tester, measuring what users experience, \qty{1.87}{\milli\second}.
\item Instants: from the intended send time to the reply's receipt at the client. Load: open loop, 1\,000 requests a
second for 100 seconds. Histogram: log-linear, 0.8\% precision, 100\,000 samples; p50 \qty{0.10}{\milli\second}, p99
\qty{1.87}{\milli\second}, p99.9 \qty{10.1}{\milli\second}. Plus machine, placement, versions, clock and repetitions.
\item A garbage collector's periodic full collection, a log rotation, a cron job, a timer interrupt storm: any periodic
work that stops the service for a few milliseconds.
\item Look for a closed-loop generator (one outstanding request per client), a sample count below the intended rate times
the duration, or a p99.9 far above a flat p99.
\item A strategy that timestamps messages when its loop reaches them rather than when they arrived.
\item Because it is inside every figure: a \qty{20}{\nano\second} clock doubles a \qty{20}{\nano\second} stage.
\item The time between the packet's arrival on the wire and the program's first read of it: the receive path and any
queueing in the socket, which software timestamps cannot see.
\item The two instants, who timestamped them, the load, the percentiles with their count, and the conditions.
\end{enumerate}
\end{solution}

\begin{solution}{iq:ll:measuring-latency:1}
Read the \omterm{def:ll:cpu-microarchitecture:tsc}{time-stamp counter} (\texttt{rdtsc}, or \texttt{rdtscp} to wait for earlier instructions) and convert with a rate
calibrated against \texttt{CLOCK\_MONOTONIC\_RAW}. It is meaningful across cores only if the counter is invariant
(constant rate in all power states, flags \texttt{constant\_tsc} and \texttt{nonstop\_tsc}) and synchronised between
sockets, which the kernel checks at boot.
\iqlookfor{TSC, invariance, calibration against a raw clock.}
\end{solution}

\begin{solution}{iq:ll:measuring-latency:2}
A generator that waits for each reply stops sending during stalls, so the requests that would have waited are never
measured. Drive load open loop, on a schedule independent of replies, and measure from each request's intended send time;
or record with the expected interval to correct.
\iqlookfor{the closed-loop mechanism and the open-loop remedy.}
\end{solution}

\begin{solution}{iq:ll:measuring-latency:3}
A log-linear histogram with fixed buckets (exact below $2^S$, then $2^{S-1}$ buckets per power of two): constant memory,
a few instructions per record, bounded relative error, mergeable across threads; read the percentile by nearest rank
over the counts.
\iqlookfor{bounded relative error and fixed memory, not stored samples.}
\end{solution}

\begin{solution}{iq:ll:measuring-latency:4}
The mean mixes the body and the tail and hides both; a few stalls move it little while costing the most. Report p50,
p99, p99.9, the maximum and the count, from a histogram, with the conditions.
\iqlookfor{percentiles and the count.}
\end{solution}

\begin{solution}{iq:ll:measuring-latency:5}
Useful to find where average time goes in a reproducible slow case (a hot function, cache misses with hardware
counters); useless for rare tail events, which it rarely samples, and for anything the sampling interrupt disturbs.
\iqlookfor{average versus tail.}
\end{solution}

\begin{solution}{iq:ll:measuring-latency:6}
Several interleaved runs of each build on the same pinned, quiet machine; the run-to-run spread of the p99; a
statistical comparison (bootstrap or permutation of runs); and a mechanism that explains the difference. Chapter~25 builds
the gate.
\iqlookfor{repetition and a noise estimate before a conclusion.}
\end{solution}
